% Part:sets-functions-relations
% Chapter: size-of-sets
% Section: zig-zag

\documentclass[../../../include/open-logic-section]{subfiles}

\begin{document}

\olfileid[tr]{sfr}{siz}{zigzag}

\olsection{Cantor'un Zikzak Yöntemi}

\begin{explain}
Daha önce bazı ``kolay'' numaralandırmaları ele aldık. Şimdi biraz daha zor
bir örneği ele alacağız. Doğal sayı ikilileri kümesini düşünün; bu kümeyi
\oliflabeldef{sfr:set:pai:sec}{\olref[set][pai]{sec}'te şöyle
tanımlamıştık:}{şöyle tanımlıyoruz:}
\[
\Nat \times \Nat = \Setabs{\tuple{n,m}}{n,m \in \Nat}
\]
Bu sıralı ikilileri şöyle bir \emph{dizi} içinde düzenleyebiliriz:
\[
\begin{array}{ c | c | c | c | c | c}
& \mathbf 0 & \mathbf 1 & \mathbf 2 & \mathbf 3 & \dots \\
\hline
\mathbf 0 & \tuple{0,0} & \tuple{0,1} & \tuple{0,2} & \tuple{0,3} & \dots \\
\hline
\mathbf 1 & \tuple{1,0} & \tuple{1,1} & \tuple{1,2} & \tuple{1,3} & \dots \\
\hline
\mathbf 2 & \tuple{2,0} & \tuple{2,1} & \tuple{2,2} & \tuple{2,3} & \dots \\
\hline
\mathbf 3 & \tuple{3,0} & \tuple{3,1} & \tuple{3,2} & \tuple{3,3} & \dots \\
\hline
\vdots & \vdots & \vdots & \vdots & \vdots & \ddots\\
\end{array}
\]
$\Nat \times \Nat$'taki her sıralı ikilinin dizide tam bir kez yer alacağı
açıktır. Özellikle $\tuple{n,m}$, $n$'inci satır ile $m$'inci sütunda yer
alacaktır. Fakat böyle bir dizinin elemanlarını ``tek boyutlu'' bir liste
olarak nasıl düzenleriz? Aşağıdaki dizideki örüntü bunu yapmanın bir yolunu
gösterir (elbette başka birçok seçenek de vardır):
\[
\begin{array}{ c | c | c | c | c | c | c}
& \mathbf 0 & \mathbf 1 & \mathbf 2 & \mathbf 3 & \mathbf 4 &\dots \\
\hline
\mathbf 0 & 0  & 1& 3 & 6& 10 &\ldots \\
\hline
\mathbf 1 &2 & 4& 7 & 11 & \dots &\ldots \\
\hline
\mathbf 2 & 5 & 8 & 12 & \ldots & \dots&\ldots \\
\hline
\mathbf 3 & 9 & 13 & \ldots & \ldots & \dots & \ldots \\
\hline
\mathbf 4 & 14 & \ldots & \ldots & \ldots & \dots & \ldots \\
\hline
\vdots & \vdots & \vdots & \vdots & \vdots&\ldots & \ddots\\
\end{array}
\]\noindent
Bu örüntüye \emph{Cantor'un zikzak yöntemi} denir. $\Nat \times \Nat$'ı
şöyle numaralandırır:
\[
\tuple{0,0}, \tuple{0,1}, \tuple{1,0}, \tuple{0,2}, \tuple{1,1},
\tuple{2,0}, \tuple{0,3}, \tuple{1,2}, \tuple{2,1}, \tuple{3,0}, \dots
\]
Bu da aşağıdaki sonucu gösterir:
\end{explain}

\begin{prop}\ollabel{natsquaredenumerable}
$\Nat \times \Nat$ !!{enumerable}dir.
\end{prop}

\begin{proof}
$f \colon \Nat \to \Nat\times\Nat$, her $k \in \Nat$ sayısını Cantor'un
zikzak dizisinde $n$'inci satır ile $m$'inci sütunun değeri $k$ olacak
biçimdeki $\tuple{n,m} \in \Nat \times \Nat$ ikilisine götürsün.
\end{proof}

\begin{explain}
Bu teknik oldukça güzel bir biçimde genelleştirilebilir. Örneğin doğal
sayıların sıralı üçlüleri kümesini, yani
\[
\Nat \times \Nat \times \Nat = \Setabs{\tuple{n,m,k}}{n,m,k \in \Nat}
\]
kümesini numaralandırmak için de onu kullanabiliriz. $\Nat \times \Nat
\times \Nat$'ı, $\Nat \times \Nat$ ile $\Nat$'ın Kartezyen çarpımı olarak
düşünürüz; yani
\[
\Nat^3 = (\Nat \times \Nat) \times \Nat =
\Setabs{\tuple{\tuple{n,m},k}}{n, m, k
  \in \Nat }.
\]
Böylece bir ekseni $\Nat$'ın, diğer ekseni de $\Nat^2$'nin
numaralandırmasıyla işaretleyerek $\Nat^3$'ü bir diziyle numaralandırabiliriz:
\[
\begin{array}{ c | c | c | c | c | c}
& \mathbf 0 & \mathbf 1 & \mathbf 2 & \mathbf 3 & \dots \\
\hline
\mathbf{\tuple{0,0}} & \tuple{0,0,0} & \tuple{0,0,1} & \tuple{0,0,2} & \tuple{0,0,3} & \dots \\
\hline
\mathbf{\tuple{0,1}} & \tuple{0,1,0} & \tuple{0,1,1} & \tuple{0,1,2} & \tuple{0,1,3} & \dots \\
\hline
\mathbf{\tuple{1,0}} & \tuple{1,0,0} & \tuple{1,0,1} & \tuple{1,0,2} & \tuple{1,0,3} & \dots \\
\hline
\mathbf{\tuple{0,2}} & \tuple{0,2,0} & \tuple{0,2,1} & \tuple{0,2,2} & \tuple{0,2,3} & \dots\\
\hline
\vdots & \vdots & \vdots & \vdots & \vdots & \ddots \\
\end{array}
\]
Dolayısıyla Cantor'un zikzak yöntemine benzer bir yöntem kullanarak
$\Nat^3$'ün bir numaralandırmasını da elde edebiliriz. Bunu sürdürüp herhangi
bir $n$ doğal sayısı için $\Nat^n$'nin numaralandırmalarını elde edebiliriz.
Öyleyse:
\end{explain}

\begin{prop}
Her $n \in \Nat$ için $\Nat^n$ !!{enumerable}dir.
\end{prop}

\begin{prob}\label{sfr:siz:zigzag:prob:posint-n}
Her $n \in \Nat$ için $(\PosInt)^n$'nin !!{enumerable} olduğunu gösterin.
\end{prob}

\begin{prob}\label{sfr:siz:zigzag:prob:posint-star}
  $(\PosInt)^*$'ın !!{enumerable} olduğunu gösterin.
  \cref{sfr:siz:zigzag:prob:posint-n}'i varsayabilirsiniz.
\end{prob}

\end{document}
